一个\emph{一阶语言}由\emph{常项}符号、\emph{函数}符号和\emph{谓词}符号组成。函数符号和常项符号都有确定的元数。例如，在\emph{算术语言}中，我们有一个常项符号~$\Obj 0$，一个一元函数符号 $\prime$，两个二元函数符号~$+$ 和 $\times$，以及一个二元谓词符号~$<$。
由\emph{变元}、常项符号和函数符号，我们构成语言的\emph{项}。由语言的项连同其谓词符号以及\emph{等词}~$\eq$，我们构成\emph{原子公式}。进而，从原子公式出发，利用逻辑联结词 $\lnot$、$\lor$、$\land$、$\lif$、$\liff$ 以及量词 $\lforall$ 和 $\lexists$，我们构成其\emph{公式}。由于我们在构成项和公式的过程中总是注意加入必要的括号，因此每个公式都有唯一的解读方式。这使得我们可以对公式结构作归纳来定义各种概念。

公式中变元的出现有时受相应量词的管辖：如果一个变元出现在某个量词的\emph{辖域}内，则它被认为是\emph{约束的}，否则是\emph{自由的}。这些概念都有归纳定义，我们也归纳地定义将项\emph{代入}公式中变元的操作。不含自由变元出现的公式称为\emph{语句}。